Las pruebas de Kruskal--Wallis evidenciaron diferencias estadísticamente significativas entre bibliotecas en los tres ejes analizados, con los siguientes resultados: Lux\_promedio (H=583.10, p=0.0000, epsilon-cuadrado=0.413, efecto grande); dBA\_prom (H=544.29, p=0.0000, epsilon-cuadrado=0.377, efecto grande); Bajada\_Mediana (H=818.01, p=0.0000, epsilon-cuadrado=0.572, efecto grande). Reportar p muy pequeño (aprox. 0.0000) no implica efecto infinito; por ello se complementó con tamaño del efecto, mostrando que las diferencias tienen magnitud sustantiva y no solo significancia por tamaño muestral.

En el análisis post-hoc (Mann--Whitney con ajuste de Holm) se identificaron contrastes puntuales robustos entre pares de bibliotecas. Los principales fueron: Lux\_promedio: Biblioteca Pública Las Ferias vs Virgilio Barco Bibloteca (p-aj=0.0000, delta=-0.996); Lux\_promedio: Biblioteca Pública del Deporte - CEFE Chapinero vs Virgilio Barco Bibloteca (p-aj=0.0000, delta=-1.000); Lux\_promedio: Gabriel Garcia Marquez - Tunal vs Virgilio Barco Bibloteca (p-aj=0.0000, delta=-0.991); Lux\_promedio: Biblioteca Julio Mario Santo Domingo vs Virgilio Barco Bibloteca (p-aj=0.0000, delta=-0.987); Lux\_promedio: Biblioteca Luis Angel Arango vs Biblioteca Pública del Deporte - CEFE Chapinero (p-aj=0.0000, delta=1.000); Lux\_promedio: Biblioteca Pública del Deporte - CEFE Chapinero vs Julio Mario Santo Domingo (p-aj=0.0000, delta=-1.000). Esto respalda que la heterogeneidad no es uniforme, sino concentrada en comparaciones específicas de sedes.

El análisis de valores atípicos mostró la siguiente incidencia: Lux\_promedio: 10.9%; dBA\_prom: 3.5%; Bajada\_Mediana: 12.0%. Este hallazgo sugiere episodios operativos extremos (picos de ruido o variabilidad de red) que deben gestionarse con estrategias de monitoreo continuo, no solo con promedios.

En modelos explicativos robustos, la ocupación mostró asociación positiva con ruido (coef=-0.361, p=0.0000), mientras que su relación con velocidad de bajada fue menor (coef=0.780, p=0.0003). En términos prácticos, la presión de usuarios parece impactar más directamente el confort acústico que la capacidad de red promedio, lo que sugiere priorizar rediseño espacial y control de carga acústica en franjas críticas.
